\documentclass[../LogosReference.tex]{subfiles}
\begin{document}

% ============================================================================
% Section: Constitutive Foundation
% ============================================================================

\section{Constitutive Foundation}\label{sec:constitutive-foundation}

The \textit{Constitutive Foundation} provides the foundational semantic structure based on exact truthmaker semantics.
Evaluation is hyperintensional, distinguishing propositions that agree on truth-value across all possible worlds but differ in their exact verification and falsification conditions.

% ------------------------------------------------------------
% Syntactic Primitives
% ------------------------------------------------------------

\subsection{Syntactic Primitives}\label{sec:syntactic-primitives}

The Constitutive Foundation interprets the following syntactic primitives:

\begin{itemize}
	\item \textbf{Variables}: $v_1, v_2, v_3, \ldots$ (ranging over states)
	\item \textbf{Individual constants}: $a, b, c, \ldots$ (0-place function symbols)
	\item \textbf{n-place function symbols}: $f, g, h, \ldots$
	\item \textbf{n-place predicates}: $F, G, H, \ldots$
	\item \textbf{Sentence letters}: $p, q, r, \ldots$ (0-place predicates)
	\item \textbf{Lambda abstraction}: $\lambda x.\metaA$ (binding variable $x$ in formula $\metaA$)
	\item \textbf{First-order identity}: $=$ (identity between terms)
	\item \textbf{Logical connectives}: $\neg, \land, \lor, \top, \bot, \equiv$
\end{itemize}

\begin{definition}[Well-Formed Formulas]\label{def:wff}
The set of \emph{well-formed formulas} is defined inductively where $F$ over $n$-place predicates, $t_i$ over terms, and $x$ over variables:
  \[
    \metaA, \metaB \Define{} F(t_1, \ldots, t_n) \mid \top \mid \bot
    \mid{} \neg\metaA \mid (\metaA \land \metaB) \mid (\metaA \lor \metaB)
    \mid{} (\metaA \equiv \metaB) \mid (\lambda x.\metaA)(t)
  \]

\end{definition}

\noindent
See \leansrc{Logos.SubTheories.Foundation.Syntax}{ConstitutiveFormula} for the Lean implementation.

\subsection{Derived Operators}\label{sec:derived-operators}

% FIX: combine the following definitions and the definition of the quantifier notation below into a single definition which sets metalinguistic conventions (syntactic sugar) for the language

\begin{definition}[Material Conditional]\label{def:material-conditional}
	\[
		\metaA \to \metaB := \neg\metaA \lor \metaB
	\]
\end{definition}

\begin{definition}[Biconditional]\label{def:biconditional}
	\[
		\metaA \leftrightarrow \metaB := (\metaA \to \metaB) \land (\metaB \to \metaA)
	\]
\end{definition}

\begin{remark}[Quantifier Notation]
First-order quantifiers are \textit{not} primitive in the Constitutive Foundation.
Generalized quantification is expressed via lambda abstraction:
\begin{itemize}
	\item Universal: $\forall x.\metaA$ abbreviates $\forall(\lambda x.\metaA)$ where $\forall$ is a second-order predicate on properties
	\item Existential: $\exists x.\metaA$ abbreviates $\exists(\lambda x.\metaA)$ where $\exists$ is a second-order predicate on properties
\end{itemize}
First-order quantification is introduced in the Dynamical Foundation layer.
\end{remark}

\noindent
See \leansrc{Logos.Dynamics.Syntax}{CoreFormula}.

% ------------------------------------------------------------
% Constitutive Frame
% ------------------------------------------------------------

\subsection{Constitutive Frame}\label{sec:constitutive-frame}

\begin{definition}\label{def:constitutive-frame}
	A \emph{constitutive frame} is a structure $\mframe = \tuple{\statespace, \parthood}$ where:
	\begin{itemize}
		\item $\statespace$ is a nonempty set of \emph{states}
		\item $\parthood$ is a partial order on $\statespace$ where $\tuple{\statespace, \parthood}$ is a complete lattice
	\end{itemize}
\end{definition}

\begin{remark}
	The constitutive frame is non-modal: neither possibility nor compatibility can be defined since this requires the task relation introduced in the Dynamical Foundation.
	Taking $S$ itself to consist of \textit{possible} states disrupts closure under fusion makes the conjunction of incompatible propositions empty, stripping the semantics of its capacity to propositions differing in subject-matter with no gain in simplicity or elegance.
  % The \textit{Dynamical Foundation} will introduce a task relation parameterized by durations which encode the possible evolutions of states.
\end{remark}

\begin{remark}
	The lattice structure provides:
	\begin{itemize}
		\item \textbf{Null state} $\nullstate$: The bottom element (fusion of the empty set)
		\item \textbf{Full state} $\fullstate$: The top element (fusion of all states)
		\item \textbf{State Fusion} $\fusion{s}{t}$: The least upper bound of $s$ and $t$
		\item \textbf{Pointwise Fusion} $\fusion{f}{g}(s_1,\ldots,s_n) \define \fusion{f(s_1,\ldots,s_n)}{g(s_1,\ldots,s_n)}$ 
	\end{itemize}
	See \leansrc{Logos.Foundation.Frame}{ConstitutiveFrame} for the implementation.
\end{remark}


% ------------------------------------------------------------
% Constitutive Model
% ------------------------------------------------------------

\subsection{Constitutive Model}\label{sec:constitutive-model}

% FIX: the definitions below need to be explained, where the leading ideas are that: objects are interpreted by states, and predicates are interpreted by functions from n states to a state. It is worth nothing that the state assigned to an object are called haecceities and do not specify every way that object is, but rather the ways that object is essentially (and so must be), which is then used to define its modal profile (the world-states in which that object exists) which may include that object being in all sorts of different ways. For instance, the state assigned to me does not include that I happen to be standing since I may also be belong to world-states in which I am sitting.

\begin{definition}\label{def:constitutive-model}
	A \emph{constitutive model} is a structure $\model = \tuple{\statespace, \parthood, \interp{\cdot}}$ where $\tuple{\statespace, \parthood}$ is a constitutive frame and $\interp{\cdot}$ is an \textit{interpretation function} satisfying:
	\begin{itemize}
		\item \textbf{n-place function symbols} $\interp{f} : (\Fin\;n \to \statespace) \to \statespace$
    \item \textbf{n-place predicates} $\verifiertype{F}, \falsifiertype{F} : ((\Fin\;n \to \statespace) \to \statespace) \to \Prop$ are the \textit{verifier function} and \textit{falsifier function} types which are required to satisfy:\footnote{The Lean implementation uses the mathematically equivalent Set-based formulation $\mathtt{Set}\;((\Fin\;n \to \statespace) \to \statespace)$ since in Lean 4, $\mathtt{Set}\;\alpha$ is defined as $\alpha \to \Prop$. The Set notation emphasizes the collectional aspect while the Prop notation emphasizes the logical aspect; both capture the same mathematical structure.}
		      \begin{itemize}
            \item $\fusion{f}{g} : |F|^\pm$ for any $f, g : |F|^\pm$.
			      \item $\Fusion(a_1, \ldots, a_n) \sqsubseteq f(a_1, \ldots, a_n)$ for any $f : |F|^\pm$ and $a_1, \ldots, a_n : \statespace$.
		      \end{itemize}
	\end{itemize}
\end{definition}

\begin{remark}\label{rem:trivial-cases}
	Constants and sentence letters are covered by the trivial cases of 0-place functions and 0-place predicates:
	\begin{itemize}
		\item \textbf{Constant symbols} $\interp{c} : \statespace$ is the type for 0-place functions.
		\item \textbf{Sentence letters} $\verifiertype{p}, \falsifiertype{p} : \statespace \to \Prop$ are the \textit{verifier function type} and \textit{falsifier function type} for 0-place predicates.
	\end{itemize}
\end{remark}

\noindent
See \leansrc{Logos.Foundation.Model}{ConstitutiveModel} for the implementation.

% ------------------------------------------------------------
% Variable Assignment
% ------------------------------------------------------------

\subsection{Variable Assignment}\label{sec:variable-assignment}

\begin{definition}\label{def:variable-assignment}
	A \emph{variable assignment} $\assignment : \mathrm{Var} \to \statespace$ maps each variable name to a state in the state space.
\end{definition}

\begin{notation}
	Greek letters ($\history, \alpha, \beta, \ldots$) are reserved for world-histories.
	The letter $\assignment$ (with subscripts $\assignment_1, \assignment_2, \ldots$) denotes variable assignments, chosen for compatibility across \LaTeX, markdown, and Lean notation.
\end{notation}

\begin{definition}[Term Algebra]\label{def:term-algebra}
The set $\Term$ of \emph{terms} is defined inductively:
\[
	t \Define x \mid c \mid f(t_1, \ldots, t_n)
\]
where $x$ ranges over variables, $c$ over individual constants (0-place function symbols), and $f$ over $n$-place function symbols.

The \emph{free variables} $\FV(t)$ of a term are defined by:
\begin{align*}
	\FV(x) &= \{x\} \\
	\FV(c) &= \emptyset \\
	\FV(f(t_1, \ldots, t_n)) &= \bigcup_{i=1}^n \FV(t_i)
\end{align*}

\emph{Substitution} $t[s/x]$ (replacing $x$ with $s$ in $t$) is defined by:
\begin{align*}
	y[s/x] &= \begin{cases} s & \text{if } y = x \\ y & \text{otherwise} \end{cases} \\
	c[s/x] &= c \\
	f(t_1, \ldots, t_n)[s/x] &= f(t_1[s/x], \ldots, t_n[s/x])
\end{align*}
\end{definition}

\noindent
See \leansrc{Logos.SubTheories.Foundation.Syntax}{Term}.

\begin{definition}\label{def:variant-assignment}
	The \textit{$x$-variant} $\assignsub{v}{x}$ of the variable assignment $\assignment$ differs at most from $\assignment$ by assigning $x$ to $v$.
\end{definition}

\begin{definition}\label{def:term-extension}
	The \emph{extension} of a term given a model $\model$ and assignment $\assignment$ is a function $\sem{\cdot}^\assignment_\model : \Term \to \statespace$ defined by:
	\begin{itemize}
		\item \textbf{Variable} $\sem{x}^\assignment_\model = \assignment(x)$
		\item \textbf{Function application} $\sem{f(t_1, \ldots, t_n)}^\assignment_\model = \interp{f}(\sem{t_1}^\assignment_\model, \ldots, \sem{t_n}^\assignment_\model)$
	\end{itemize}
\end{definition}

\noindent
See \leansrc{Logos.Foundation.Assignment}{Assignment}.

% ------------------------------------------------------------
% Verification and Falsification Clauses
% ------------------------------------------------------------

\subsection{Verification and Falsification}\label{sec:verification-falsification}

The \emph{verification} ($\verifies$) and \emph{falsification} ($\falsifies$) of a formula $\metaA$ is defined recursively relative to a model $\model$, assignment $\assignment$, and state $s$.

\subsubsection{Atomic Formulas}\label{sec:atomic-formulas}

\begin{definition}[Atomic Verification and Falsification]\label{def:atomic-verification}
	\begin{align}
		\model, \assignment, s \verifies F(t_1, \ldots, t_n)  & \iff \exists f \in \verifiertype{F} : s = f(\lambda i : \Fin\;n, \sem{t_i}^\assignment_\model)  \\
		\model, \assignment, s \falsifies F(t_1, \ldots, t_n) & \iff \exists f \in \falsifiertype{F} : s = f(\lambda i : \Fin\;n, \sem{t_i}^\assignment_\model)
	\end{align}
\end{definition}

\subsubsection{Lambda Abstraction}\label{sec:lambda-abstraction}

\begin{definition}[Lambda Verification and Falsification]\label{def:lambda-verification}
	\begin{align}
		\model, \assignment, s \verifies (\lambda x.\metaA)(t)  & \iff \model, \assignsub{\sem{t}^\assignment_\model}{x}, s \verifies \metaA  \\
		\model, \assignment, s \falsifies (\lambda x.\metaA)(t) & \iff \model, \assignsub{\sem{t}^\assignment_\model}{x}, s \falsifies \metaA
	\end{align}
\end{definition}

\subsubsection{Negation}\label{sec:negation}

\begin{definition}[Negation Verification and Falsification]\label{def:negation-verification}
	\begin{align}
		\model, \assignment, s \verifies \neg\metaA  & \iff \model, \assignment, s \falsifies \metaA \\
		\model, \assignment, s \falsifies \neg\metaA & \iff \model, \assignment, s \verifies \metaA
	\end{align}
\end{definition}

\begin{remark}
Negation exchanges verifiers and falsifiers: what makes $\metaA$ true makes $\neg\metaA$ false, and vice versa.
This maintains bilateral symmetry and ensures that a state wholly relevant to $\metaA$ is equally relevant to $\neg\metaA$.
\end{remark}

\subsubsection{Conjunction}\label{sec:conjunction}

\begin{definition}[Conjunction Verification and Falsification]\label{def:conjunction-verification}
	\begin{align}
		\model, \assignment, s \verifies \metaA \land \metaB \iff
		 & s = \fusion{t}{u} \text{ for some } t, u \text{ where } \model, \assignment, t \verifies \metaA \notag                                           \\
		 & \text{and } \model, \assignment, u \verifies \metaB                                                                                              \\
		\model, \assignment, s \falsifies \metaA \land \metaB \iff
		 & \model, \assignment, s \falsifies \metaA, \text{ or } \model, \assignment, s \falsifies \metaB, \text{ or } s = \fusion{t}{u} \text{ for} \notag \\
		 & \text{some } t, u \text{ where } \model, \assignment, t \falsifies \metaA \text{ and } \model, \assignment, u \falsifies \metaB
	\end{align}
\end{definition}

\begin{remark}
A conjunction $\metaA \land \metaB$ is verified by a state that is the \textit{fusion} of a verifier for $\metaA$ and a verifier for $\metaB$.
If state $t$ verifies ``Alice is thinking'' and state $u$ verifies ``Alice is writing,'' then $\fusion{t}{u}$ verifies ``Alice is thinking and writing.''
Crucially, $\fusion{t}{u}$ does \textit{not} verify ``Alice is thinking'' alone since $u$ is irrelevant to that claim. The non-monotonicity of the verification and falsification relations is characteristic of exact semantics.

The falsification clause is \textit{inclusive}: a conjunction is falsified by anything that falsifies either conjunct, or by a fusion of falsifiers for both conjuncts.
This ensures that $\metaA \land \metaA$ has the same falsifiers as $\metaA$ so that $\metaA \land \metaA \equiv \metaA$ comes out valid, providing a natural idempotence law.
\end{remark}

\subsubsection{Disjunction}\label{sec:disjunction}

\begin{definition}[Disjunction Verification and Falsification]\label{def:disjunction-verification}
	\begin{align}
		\model, \assignment, s \verifies \metaA \lor \metaB \iff
		 & \model, \assignment, s \verifies \metaA, \text{ or } \model, \assignment, s \verifies \metaB, \text{ or } s = \fusion{t}{u} \text{ for} \notag \\
		 & \text{some } t, u \text{ where } \model, \assignment, t \verifies \metaA \text{ and } \model, \assignment, u \verifies \metaB                  \\
		\model, \assignment, s \falsifies \metaA \lor \metaB \iff
		 & s = \fusion{t}{u} \text{ for some } t, u \text{ where } \model, \assignment, t \falsifies \metaA \notag                                        \\
		 & \text{and } \model, \assignment, u \falsifies \metaB
	\end{align}
\end{definition}

\begin{remark}[Disjunction Semantics]
Disjunction is dual to conjunction: verification and falsification swap roles.
A disjunction $\metaA \lor \metaB$ is falsified by a fusion of falsifiers for each disjunct---the state must be wholly relevant to falsifying both.
The verification clause is inclusive, mirroring the falsification clause for conjunction, making $\metaA \lor \metaA \equiv \metaA$ valid.
\end{remark}

\subsubsection{Top and Bottom}\label{sec:top-bottom}

\begin{definition}[Top and Bottom Verification and Falsification]\label{def:top-bottom-verification}
	\begin{align}
		 & \model, \assignment, s \verifies \top \text{ for all } s \in \statespace                 \\
		 & \model, \assignment, s \falsifies \top \iff s = \fullstate \text{ (\textit{full state})} \\
		 & \model, \assignment, s \nverifies \bot \text{ for all } s                                \\
		 & \model, \assignment, s \falsifies \bot \iff s = \nullstate \text{ (\textit{null state})}
	\end{align}
\end{definition}

\subsubsection{Propositional Identity}\label{sec:propositional-identity}

\begin{definition}[Propositional Identity Verification and Falsification]\label{def:propid-verification}
	\begin{align}
		\model, \assignment, s \verifies \metaA \equiv \metaB  \iff
		 & s = \nullstate \text{ and for all } t \text{ both:} \notag                                                                     \\
		 & (1)~ \model, \assignment, t \verifies \metaA \text{ just in case } \model, \assignment, t \verifies \metaB; \text{ and} \notag \\
		 & (2)~ \model, \assignment, t \falsifies \metaA \text{ just in case } \model, \assignment, t \falsifies \metaB                   \\
		\model, \assignment, s \falsifies \metaA \equiv \metaB \iff
		 & \text{ at least one of the following:} \notag                                                                                                         \\
		 & (1)~ \model, \assignment, s \verifies \metaA \text{ and } \model, \assignment, s \nverifies \metaB; \notag                     \\
		 & (2)~ \model, \assignment, s \nverifies \metaA \text{ and } \model, \assignment, s \verifies \metaB; \notag                     \\
		 & (3)~ \model, \assignment, s \falsifies \metaA \text{ and } \model, \assignment, s \nfalsifies \metaB; \text{ or} \notag        \\
		 & (4)~ \model, \assignment, s \nfalsifies \metaA \text{ and } \model, \assignment, s \falsifies \metaB.
	\end{align}
\end{definition}

\noindent
See \leansrc{Logos.Foundation.Semantics}{verifies} for the implementation.

% ------------------------------------------------------------
% Essence and Ground
% ------------------------------------------------------------

\subsection{Essence and Ground}\label{sec:essence-ground}

These constitutive relations can be defined via propositional identity:

\begin{definition}[Essence and Ground]\label{def:essence-ground}
	\begin{align}
		\metaA \essence \metaB & \define \metaA \land \metaB \equiv \metaB &  & \text{``$\metaA$ is essential to $\metaB$'' (conjunctive part)} \\
		\metaA \ground \metaB  & \define \metaA \lor \metaB \equiv \metaB  &  & \text{``$\metaA$ grounds $\metaB$'' (disjunctive part)}
	\end{align}
\end{definition}

\begin{definition}[Reduction]\label{def:reduction}
  \emph{Reduction} combines essence and ground:
  \[
    \metaA \Rightarrow \metaB \define (\metaA \essence \metaB) \land (\metaA \ground \metaB)
  \]
  Reduction expresses that $\metaA$ is both essential to (necessary for) and grounds (sufficient for) $\metaB$.
  Examples include $\metaA \Rightarrow \metaA \lor (\metaA \land \metaB)$ and $\metaA \Rightarrow \metaA \land (\metaA \lor \metaB)$. 
\end{definition}

\begin{remark}
	Essence and ground are interrelated by negation:
	\begin{itemize}
		\item $\metaA \essence \metaB \iff \neg\metaA \ground \neg\metaB$
		\item $\metaA \ground \metaB \iff \neg\metaA \essence \neg\metaB$
	\end{itemize}
\end{remark}

\noindent
See \leansrc{Logos.Foundation.Relations}{essence}.

% ------------------------------------------------------------
% Bilateral Propositions
% ------------------------------------------------------------

\subsection{Bilateral Propositions}\label{sec:bilateral-propositions}

\begin{definition}[Bilateral Proposition]\label{def:bilateral-proposition}
	A \emph{bilateral proposition} $P : \props$ is an ordered pair $\tuple{V, F}$ where $V$ and $F$ are subsets of $\statespace$ closed under fusion:
	\begin{itemize}
		\item \textbf{Fusion closure for verifiers}: $\fusion{s}{t} \in V$ for any $s, t \in V$
		\item \textbf{Fusion closure for falsifiers}: $\fusion{s}{t} \in F$ for any $s, t \in F$
	\end{itemize}
\end{definition}

\begin{definition}[Bilattice]\label{def:bilattice}
A \emph{bilattice} is a structure $B = \tuple{\props, \essence, \ground, \neg}$ where:
\begin{itemize}
	\item $\tuple{\props, \essence}$ is a complete lattice (the \emph{essence} order)
	\item $\tuple{\props, \ground}$ is a complete lattice (the \emph{ground} order)
	\item $\neg : \props \to \props$ is an \emph{involutive negation}: $\neg\neg P = P$
	\item Negation exchange: $P \essence Q \leftrightarrow \neg P \ground \neg Q$ and $P \ground Q \leftrightarrow \neg P \essence \neg Q$.
\end{itemize}
A bilattice is \emph{interlaced} when its orders interact uniformly:
$(P \star R) \circ (Q \star R)$ whenever $P \circ Q$, for any lattice operation $\star$ and either order $\circ$.
\end{definition}

\begin{remark}\label{rem:bilattice-orders}
  The essence order $P \essence Q$ captures a constitutive explanatory reading of ``necessary for'' whereby: (1) every verifier for $Q$ has a part that verifies $P$; (2) any fusion of a verifier for $P$ together with a verifier for $Q$ also verifies $Q$; and (3) every falsifier for $P$ also falsifies $Q$.
  Intuitively, $P$ being the case is \textit{part of what it is for} $Q$ to be the case.

  The ground order $P \ground Q$ captures a constitutive explanatory reading of ``sufficient for'' whereby: (1) every verifier for $P$ also verifies $Q$; (2) every falsifier for $Q$ has a part that falsifies $P$; and (3) any fusion of a falsifier for $P$ together with a falsifier for $Q$ also falsifies $Q$.
  Intuitively, $P$ being the case \textit{determines} that it is the case that $Q$.
\end{remark}

\begin{remark}\label{rem:bilattice-structure}
  The bilateral propositions form a \emph{non-interlaced} bilattice:
  \begin{itemize}
    \item Conjunction is the least upper bound with respect to $\essence$
    \item Disjunction is the least upper bound with respect to $\ground$
    \item The two orders need not interact uniformly (hence non-interlaced)
  \end{itemize}
  The failure of interlacing reflects the hyperintensionality of the semantics:
  giving up Boolean identities allows the two orders to diverge.
\end{remark}

\begin{definition}[Propositional Operations]\label{def:propositional-operations}
	For sets of states $X, Y \subseteq \statespace$ and bilateral propositions $P = \tuple{P^+, P^-}$ and $Q = \tuple{Q^+, Q^-}$, we may define the operations on their verifier and falsifier sets as follows:
	\begin{itemize}
		\item \textbf{Pairwise Fusion}: $X \otimes Y \define \set{\fusion{s}{t} : s \in X, t \in Y}$
		\item \textbf{Closed Fusion}: $X \oplus Y \define X \cup Y \cup (X \otimes Y)$
		\item \textbf{Conjunction}: $P \land Q \define \tuple{P^+ \otimes Q^+, P^- \oplus Q^-}$
		\item \textbf{Disjunction}: $P \lor Q \define \tuple{P^+ \oplus Q^+, P^- \otimes Q^-}$
		\item \textbf{Negation}: $\neg P \define \tuple{P^-, P^+}$
	\end{itemize}
\end{definition}

\noindent
See \leansrc{Logos.Foundation.Proposition}{BilateralProp}.

% ------------------------------------------------------------
% Logical Consequence
% ------------------------------------------------------------

\subsection{Logical Consequence}\label{sec:logical-consequence-foundation}

% FIX: the identity sentences need to be precisely defined in both the lean source code and here, where this inductive definition is in terms of the well-formed sentences but takes identity sentences to be atomic

\begin{definition}[Logical Consequence]\label{def:consequence-foundation}
	The \textit{Constitutive Foundation} of the Logos is restricted to propositional identity sentences of the form $\metaA \equiv \metaB$ where $\metaA$ and $\metaB$ are atomic sentences or built from $\neg, \land, \lor, \equiv$.
	\begin{align*}
		\context \trueat \metaA \iff
      & \text{for any model } \model \text{ and assignment } \assignment : \\
      & \text{if } \model, \assignment, \nullstate \trueat \metaB \text{ for all } \metaB \in \context, \text{ then } \model, \assignment, \nullstate \trueat \metaA
	\end{align*}
\end{definition}

\begin{remark}[Identity and Contingency]
	The \textit{Constitutive Foundation} does not provide the semantic resources to evaluate non-identity sentences which depend on contingent states of affairs rather than purely structural relations in state space.
	The \textit{Dynamical Foundation} encodes contingencies by including a world-history, time, and list of temporal references in addition to a variable assignment in the contexts at which sentences are evaluated.
\end{remark}

\end{document}
